% Copyright (c) 2026 Selina Cheggour. Tous droits réservés. Voir le fichier LICENSE.
% Contexte, vue d'ensemble de la chaîne, mise en route sur Kaggle, règles de confidentialité.

\section{Le problème en une page}

Une entreprise agroalimentaire reçoit chaque mois des centaines de plaintes de consommateurs, dans
plusieurs pays et plusieurs langues. Chaque plainte est un court texte libre, le \textbf{verbatim},
accompagné d'un numéro de cas, de la marque (numérotée), du segment d'activité et de la période.
\textbf{Personne n'a classé ces textes.}

\begin{tcolorbox}[colback=notebg,colframe=sectioncolor,arc=3pt,halign=center,boxrule=1.5pt]
{\Large\faSearch\quad\bfseries Quand plusieurs consommateurs décrivent le même problème,}\\[2pt]
{\Large\bfseries peut-on les \textcolor{accentblue}{regrouper} et voir si ce problème se
\textcolor{accentblue}{concentre} sur une marque et une période ?}
\end{tcolorbox}

\begin{enfant}
\textbf{Comme le courrier d'une mairie.} Si dix habitants écrivent la même semaine que l'eau du robinet
a un goût bizarre, c'est sans doute une canalisation, pas dix coïncidences. Mais si la moitié de la
ville se plaint du froid en janvier, ce n'est pas une panne : c'est l'hiver. Et si mille lettres
identiques arrivent le même jour, c'est une pétition. Le TP apprend à faire la différence.
\end{enfant}

\subsection{Trois phénomènes qui se ressemblent, et qu'il faut séparer}

\begin{center}
\begin{tikzpicture}[x=0.32cm, y=0.13cm]
% trois mini-histogrammes : rouge = la marque concernée, beige = les autres marques
\foreach \nom/\dx/\titre/\coul in {inc/0/Incident/accentblue, sai/17/Saison/ambre, cam/34/Campagne/cacao} {
  \node[font=\bfseries, color=\coul] at (\dx+6.5,30) {\titre};
  \draw[slate!60] (\dx-0.5,0) -- (\dx+13.5,0);
}
% incident : une seule marque explose sur une période
\foreach \x/\a/\m in {0/3/0,1/2/1,2/4/0,3/3/1,4/2/0,5/3/0,6/3/16,7/4/2,8/2/0,9/3/1,10/2/0,11/3/0,12/2/0,13/3/1} {
  \fill[palblue] (\x-0.4,0) rectangle (\x+0.4,\a); \fill[accentblue] (\x-0.4,\a) rectangle (\x+0.4,\a+\m);}
% saison : toutes les marques montent ensemble
\foreach \x/\a/\m in {0/3/1,1/3/0,2/4/1,3/5/1,4/9/3,5/14/3,6/16/4,7/13/3,8/7/2,9/4/1,10/3/0,11/3/1,12/2/0,13/3/0} {
  \fill[palblue] (17+\x-0.4,0) rectangle (17+\x+0.4,\a); \fill[accentblue] (17+\x-0.4,\a) rectangle (17+\x+0.4,\a+\m);}
% campagne : un gros paquet de messages presque identiques
\foreach \x/\a/\m in {0/1/0,1/1/0,2/2/0,3/1/0,4/1/0,5/12/6,6/14/8,7/9/4,8/1/0,9/2/0,10/1/0,11/1/0,12/1/0,13/2/0} {
  \fill[palblue] (34+\x-0.4,0) rectangle (34+\x+0.4,\a); \fill[accentblue] (34+\x-0.4,\a) rectangle (34+\x+0.4,\a+\m);}
\node[font=\footnotesize, text width=4.9cm, align=center] at (6.5,-6) {une \textbf{seule marque} explose sur une ou deux périodes : un lot, une usine, une recette};
\node[font=\footnotesize, text width=4.9cm, align=center] at (23.5,-6) {\textbf{toutes les marques} montent ensemble : le chocolat fond l'été};
\node[font=\footnotesize, text width=4.9cm, align=center] at (40.5,-6) {des centaines de messages \textbf{presque identiques} : une mobilisation, pas un défaut};
\fill[accentblue] (12,-13) rectangle (12.8,-11.5); \node[right, font=\footnotesize] at (12.9,-12.2) {la marque signalée};
\fill[palblue] (24,-13) rectangle (24.8,-11.5); \node[right, font=\footnotesize] at (24.9,-12.2) {les autres marques};
\end{tikzpicture}
\end{center}
{\small\color{slate}\textit{Schéma inventé : une barre par période, pour un même problème.}}

\begin{tcbraster}[raster columns=3, raster equal height, raster column skip=6pt, raster every box/.style={unbreakable}]
\begin{kpibox}[title={\faExclamationTriangle\ Incident candidat}, unbreakable]
\small un problème en \textbf{excès} sur une case (marque, période), et cette marque porte au moins la
moitié du pic. \textbf{C'est ce qu'on cherche.}
\end{kpibox}
\begin{kpibox}[title={\faSun\ Pic partagé}, unbreakable]
\small un problème en excès sur une période, mais réparti entre les marques : une saison, un phénomène
général.
\end{kpibox}
\begin{kpibox}[title={\faBullhorn\ Campagne}, unbreakable]
\small un groupe dont la plupart des messages ont une jumelle \textbf{presque identique}.
\end{kpibox}
\end{tcbraster}

\begin{piege}
\textbf{Il n'y a pas d'étiquette.} On ne prédit rien : on \textbf{regroupe}. C'est de l'apprentissage
non supervisé, et toute la difficulté est de montrer que les groupes trouvés veulent dire quelque chose.
Le notebook travaille sur la \textbf{base réelle}. Les figures de cette fiche viennent d'une
\textbf{base fictive} de même forme, pour montrer à quoi ressemblent les résultats sans montrer de données
confidentielles.
\end{piege}

%%%%%%%%%%%%%%%%%%%%%%%%%%%%%%%%%%%%%%%%%%%%%%%%%%%%%%%%%%%%%%%%%%%%%%%%%%%%%%%%%%%%%%%%%%%%%%%%%%%%%%%%%
\newpage
\section{Deux grandes parties, dans cet ordre}

Le travail se fait en deux temps. D'abord la \textbf{chaîne principale} : un seul outil, de la lecture
du fichier jusqu'au tableau de bord. Ensuite seulement, une fois qu'elle marche, la \textbf{chaîne de
comparaison} : la même chaîne, rejouée avec huit outils.

\begin{center}
\begin{tikzpicture}[
  p/.style={draw=#1, fill=#1!10, rounded corners=6pt, minimum width=15.6cm, minimum height=1.3cm,
    text width=15.2cm, align=left, font=\small, inner sep=5pt},
  n/.style={circle, fill=#1, text=white, font=\bfseries, minimum size=0.9cm},
  gp/.style={fill=#1, text=white, rounded corners=4pt, font=\bfseries, minimum width=15.6cm, minimum height=0.8cm,
    text width=15.2cm, align=left, inner sep=5pt}]
\node[gp=sectioncolor] (g1) {\faRoute\quad Première grande partie : la chaîne principale \hfill{\normalfont\small un outil, Qwen3-Embedding-0.6B}};
\node[p=slate, below=0.25cm of g1] (p0) {\hspace{1.3cm}\parbox{13.7cm}{\textbf{Mise en route.} Carte graphique, internet, la base réelle attachée, outils et contrôles \texttt{verifier}.}};
\node[p=sectioncolor, below=0.25cm of p0] (p1) {\hspace{1.3cm}\parbox{13.7cm}{\textbf{Lire et contrôler.} Trouver l'en-tête tout seul, nettoyer les verbatims, compter, détecter la langue, chercher les données personnelles oubliées.}};
\node[p=darkblue, below=0.25cm of p1] (p2) {\hspace{1.3cm}\parbox{13.7cm}{\textbf{Transformer en vecteurs.} TF-IDF puis plongements multilingues ; retirer la \og signature \fg{} de chaque langue par centrage.}};
\node[p=accentblue, below=0.25cm of p2] (p3) {\hspace{1.3cm}\parbox{13.7cm}{\textbf{Regrouper.} HDBSCAN, fusion des groupes jumeaux d'une langue à l'autre, réglage automatique, mots et définitions de chaque problème, stabilité.}};
\node[p=lightblue, below=0.25cm of p3] (p4) {\hspace{1.3cm}\parbox{13.7cm}{\textbf{Repérer les incidents.} Plaintes observées contre attendues (loi de Poisson), correction des tests multiples, part de la marque dans le pic ; incident, saison ou campagne.}};
\node[p=cacao, below=0.25cm of p4] (p6) {\hspace{1.3cm}\parbox{13.7cm}{\textbf{Exporter.} \texttt{resultats.js} pour le tableau de bord, contrôlé : ni verbatim ni numéro de cas.}};
\node[draw=accentblue, fill=white, line width=1.2pt, rounded corners=12pt, below=0.35cm of p6, font=\small\bfseries,
  text=accentblue, inner sep=6pt] (porte) {\faCheckCircle\quad la chaîne principale est validée : ses vérifications lues, \texttt{resultats.js} écrit};
\node[gp=bordeaux, below=0.35cm of porte] (g2) {\faBalanceScale\quad Deuxième grande partie : la chaîne de comparaison \hfill{\normalfont\small huit outils}};
\node[p=bordeaux, below=0.25cm of g2] (p7) {\hspace{1.3cm}\parbox{13.7cm}{\textbf{Comparer les outils.} La même chaîne, du centrage aux alertes, avec huit représentations, du comptage de mots à un transformer de 4 milliards de paramètres : temps, mémoire, coût, qualité. Sortie : \texttt{comparaison.js}.}};
\node[p=slate, below=0.35cm of p7] (p8) {\hspace{1.3cm}\parbox{13.7cm}{\textbf{Pour finir : ce qui peut sortir de Kaggle.} Des chiffres seulement ; aucun texte ne s'affiche ni ne sort.}};
\foreach \k/\c in {0/slate,1/sectioncolor,2/darkblue,3/accentblue,4/lightblue,6/cacao,7/bordeaux,8/slate}
  \node[n=\c] at ($(p\k.west)+(0.75,0)$) {\k};
\draw[-{Stealth[length=3mm]}, very thick, color=accentblue] (p6) -- (porte);
\draw[-{Stealth[length=3mm]}, very thick, color=accentblue] (porte) -- (g2);
\end{tikzpicture}
\end{center}

\begin{enfant}
\textbf{Comme apprendre une recette avant d'acheter des casseroles.} On réussit d'abord le plat avec une
seule casserole (la chaîne principale). Ensuite seulement, on refait exactement la même recette avec huit
casseroles pour voir laquelle cuit le mieux, le plus vite et pour le moins cher (la chaîne de comparaison).
Comparer des casseroles sur une recette qu'on rate encore ne dirait rien.
\end{enfant}

%%%%%%%%%%%%%%%%%%%%%%%%%%%%%%%%%%%%%%%%%%%%%%%%%%%%%%%%%%%%%%%%%%%%%%%%%%%%%%%%%%%%%%%%%%%%%%%%%%%%%%%%%
\newpage
\section{Mise en route sur Kaggle}

\begin{bloc}[darkblue]{\faCalendarCheck\quad À faire AVANT la séance}
\begin{tikzpicture}[
  etape/.style={draw=darkblue,fill=white,rounded corners=4pt,minimum width=3.6cm,minimum height=1.5cm,
    align=center,font=\small,text width=3.4cm},
  arr/.style={-{Stealth[length=3mm]},very thick,color=darkblue}]
\node[etape] (a) {{\Large\faUserPlus}\\Créer un compte\\\textbf{kaggle.com}};
\node[etape,right=0.35cm of a] (b) {{\Large\faMobile*}\\\textbf{Vérifier ton téléphone}\\(sinon : ni GPU ni internet)};
\node[etape,right=0.35cm of b] (c) {{\Large\faFileImport}\\Importer le notebook\\du TP (\textbf{File, Import})};
\node[etape,right=0.35cm of c] (d) {{\Large\faDatabase}\\\textbf{Add Input} : le jeu\\de la base réelle};
\draw[arr] (a) -- (b); \draw[arr] (b) -- (c); \draw[arr] (c) -- (d);
\end{tikzpicture}
\end{bloc}

\begin{bloc}{\faPowerOff\quad Allumer le GPU et internet}
\begin{tikzpicture}[
  num/.style={circle,fill=sectioncolor,text=white,font=\bfseries\Large,minimum size=0.9cm},
  lab/.style={align=center,font=\small,text width=3.6cm},
  arr/.style={-{Stealth[length=3mm]},very thick,color=sectioncolor}]
\node[num] (n1) {1}; \node[num,right=3.0cm of n1] (n2) {2};
\node[num,right=3.0cm of n2] (n3) {3}; \node[num,right=3.0cm of n3] (n4) {4};
\node[lab,below=2pt of n1] {Panneau de droite :\\\textbf{Session options}};
\node[lab,below=2pt of n2] {\textbf{Accelerator}\\\faArrowRight\ \textbf{GPU T4}};
\node[lab,below=2pt of n3] {\textbf{Internet}\\\faToggleOn\ \textbf{On}};
\node[lab,below=2pt of n4] {Exécuter les cellules\\de la \textbf{partie 0}};
\draw[arr] (n1) -- (n2); \draw[arr] (n2) -- (n3); \draw[arr] (n3) -- (n4);
\end{tikzpicture}
\end{bloc}

\begin{tcbraster}[raster columns=2, raster equal height, raster column skip=8pt, raster every box/.style={unbreakable}]
\begin{enfant}[unbreakable]
\textbf{Internet sert à télécharger les modèles} (des fichiers de plusieurs centaines de mégaoctets
à quelques gigaoctets), jamais à envoyer les plaintes : tous les calculs se font dans le notebook.
\end{enfant}
\begin{enfant}[unbreakable]
\texttt{[À revoir]} = la marge d'une copie corrigée : une remarque et un \textbf{indice}, au lieu d'une
erreur qui arrête tout. \textbf{Lis l'indice} avant de passer à la suite.
\end{enfant}
\end{tcbraster}

\clearpage
\section{Les règles de confidentialité}

Le notebook tourne sur des plaintes réelles, couvertes par un accord de confidentialité que tu as signé.
Les règles ci-dessous s'appliquent donc à chaque exécution.

\begin{center}
\resizebox{\linewidth}{!}{%
\begin{tikzpicture}[
  z/.style={draw=#1, fill=#1!8, rounded corners=8pt, minimum height=3.5cm, text width=4.1cm, align=center, font=\small},
  f/.style={-{Stealth[length=3mm]}, very thick, color=#1}]
\node[z=darkblue] (k) {{\huge\faLock}\\[3pt]\textbf{Kaggle, en privé}\\[3pt]les verbatims, les numéros de cas, les vecteurs, les mots des problèmes\\[3pt]\textit{tout reste ici}};
\node[z=accentblue, right=1.7cm of k] (e) {{\huge\faShield*}\\[3pt]\textbf{L'export contrôlé}\\[3pt]\texttt{resultats.js}, \texttt{comparaison.js} : des chiffres seulement\\[3pt]\textit{refusé s'il contient un numéro de cas ou un verbatim}};
\node[z=lightblue, right=1.7cm of e] (t) {{\huge\faChartPie}\\[3pt]\textbf{Le tableau de bord}\\[3pt]sur ta machine, sans internet\\[3pt]\textit{les problèmes y sont numérotés}};
\draw[f=darkblue] (k) -- node[above, font=\footnotesize] {calcul} (e);
\draw[f=accentblue] (e) -- node[above, font=\footnotesize] {téléchargement} (t);
\end{tikzpicture}}
\end{center}

\begin{tikzpicture}[
  rule/.style={draw=#1,fill=#1!10,rounded corners=5pt,minimum height=1.75cm,
    text width=7.55cm,align=left,font=\normalsize,inner sep=8pt},
  ico/.style={font=\Huge,color=#1}]
\node[rule=sectioncolor] (r1) at (0,0) {\hspace{1.4cm}\parbox{5.95cm}{\textbf{Ne recopie aucun verbatim}\\
  ni dans un rapport, ni dans un message, ni dans une capture d'écran.}};
\node[ico=sectioncolor] at ($(r1.west)+(0.8,0)$) {\faCopy};
\node[rule=accentblue,right=0.3cm of r1] (r2) {\hspace{1.4cm}\parbox{5.95cm}{\textbf{Aucun service en ligne}\\
  pour lire ou traduire les plaintes : pas de chatbot, pas de traducteur web.}};
\node[ico=accentblue] at ($(r2.west)+(0.8,0)$) {\faCloud};
\node[rule=darkblue,below=0.3cm of r1] (r3) {\hspace{1.4cm}\parbox{5.95cm}{\textbf{Notebook privé.}\\
  Jamais \og Public \fg{}, jamais partagé hors de l'équipe.}};
\node[ico=darkblue] at ($(r3.west)+(0.8,0)$) {\faUserSecret};
\node[rule=lightblue,below=0.3cm of r2] (r4) {\hspace{1.4cm}\parbox{5.95cm}{\textbf{Ne télécharge que l'export}\\
  (\texttt{tableau\_de\_bord/*.js}), jamais les vecteurs \texttt{plongements\_*.npy} : on sait parfois
  retrouver une partie d'un texte à partir de son vecteur.}};
\node[ico=lightblue] at ($(r4.west)+(0.8,0)$) {\faDownload};
\end{tikzpicture}
